\section{Selected Structural Comparison}
\label{sec:selected-comparison}

The preceding sections put two main structural tests on the same ambient set \(\Sfour\): exact poset-cone status and subgroup-row status. The comparison in this section uses selected witnesses only. It is not a classification of all selected data rows and not a classification of all subsets of \(\Sfour\).

In \cref{tab:selected-witnesses}, the induced-cosets column is bookkeeping induced by subgroup status under \cref{def:composition-convention}. A ``yes'' in that column means that the row is a subgroup \(H\le\Sfour\), so its left and right cosets partition \(\Sfour\). It is not a separate geometric tiling statement and is not used as an independent structural axis.

\begin{table}[htbp]
\centering
\small
\begin{tabular}{@{}p{0.18\textwidth}ccc p{0.30\textwidth}@{}}
\toprule
Witness & Exact & Subgroup & Induced cosets & Role \\
\midrule
\(\{\word{1234}\}\) & yes & yes & yes & overlap at the identity row \\
\(\Sfour\) & yes & yes & yes & overlap at the full group row \\
N-poset row & yes & no & no & exact poset cone outside subgroup rows \\
\(\Dfour\) row & no & yes & yes & subgroup row outside exact poset cones \\
\(\Vfour\) row & no & yes & yes & subgroup row outside exact poset cones \\
\bottomrule
\end{tabular}
\caption{Selected witnesses used for \cref{prop:selected-structural-comparison}. The induced-cosets column is subgroup bookkeeping. The table compares only the displayed rows and is not a full data-package classification.}
\label{tab:selected-witnesses}
\end{table}

\begin{theorem}[Selected witnesses separate poset-cone and subgroup-row status]
\label{prop:selected-structural-comparison}
The displayed rows in \cref{tab:selected-witnesses} witness that exact poset-cone status and subgroup-row status are not equivalent predicates on subsets of \(\Sfour\).
\end{theorem}

\begin{proof}
By \cref{prop:n-poset-witness}, the N-poset row is an exact poset cone and has five elements. Since \(5\nmid 24\), it is not a subgroup row of \(\Sfour\). Thus exact poset-cone status does not imply subgroup-row status.

Conversely, by \cref{cor:d4-v4-not-poset-cones,cor:d4-v4-subgroup-cosets}, the fixed \(\Dfour\) and \(\Vfour\) rows are subgroup rows, but neither is an exact poset cone. Hence subgroup-row status does not imply exact poset-cone status.

The two overlap rows in \cref{tab:selected-witnesses} record that the fields are not disjoint. The singleton \(\{\word{1234}\}\) is \(\Lin(P)\) for the chain \(1<2<3<4\), and \(\word{1234}\) is the identity subgroup under the group convention. The full row \(\Sfour\) is \(\Lin(P)\) for the antichain on \([4]\), and it is also the full subgroup of itself. These overlap rows are not needed to prove non-equivalence, but they prevent the stronger and false reading that poset-cone and subgroup rows are disjoint classes.
\end{proof}

\begin{remark}[Boundary of the comparison]
\label{rem:comparison-boundary}
\Cref{prop:selected-structural-comparison} separates exact poset-cone status from subgroup-row status. It does not prove a separate pairwise non-equivalence theorem between subgroup-row status and subgroup-coset partition status. In this paper, the induced coset entries for the displayed subgroup rows are the standard left and right coset partitions induced by subgroup status. A full selected-data comparison theorem would require a new row-by-row statement and a separate certificate.
\end{remark}